\chapter{Cosmic Impedance: The Mechanics of Dark Matter}
\label{app:cosmic_impedance}

\noindent\textbf{Subject:} Galaxy-scale acceleration generated by the coherent reactive response of the Vacuum Interstitial Condensate.

In standard cosmology, galactic mass discrepancies are modeled primarily through an additional non-baryonic matter component. Tension Field Theory attributes the corresponding excess acceleration to the mechanical response of the Vacuum Interstitial Condensate within the active FCC Soft Floor.

A baryonic tension field phase-loads the surrounding condensate and establishes a correlated restorative response. In the low-acceleration regime, this response contributes to orbital acceleration and appears, under an inverse-square interpretation, as additional dynamical mass. The galactic response is reactive rather than dissipative: the condensate stores and returns deformation energy without producing continuous orbital drag.

\section{Derivation of the Lattice Impedance Acceleration}

The active FCC Soft Floor possesses the baseline surface inertia
\begin{equation}
    \mu_{\mathrm{D}},
\end{equation}
with dimensions
\begin{equation}
    [\mu_{\mathrm{D}}] = \mathrm{kg/m^2}.
\end{equation}

The corresponding lattice surface-tension scale is
\begin{equation}
    \tau_{\mathrm{D}} = \mu_{\mathrm{D}}c_{\mathrm{E}}^2,
\end{equation}
where $c_{\mathrm{E}}$ is the characteristic elastic propagation speed of the active vacuum. The dimensions of $\tau_{\mathrm{D}}$ are
\begin{equation}
    [\tau_{\mathrm{D}}] = \mathrm{N/m}.
\end{equation}

Unlike the inverse-square field generated by a localized source, the background lattice impedance is associated with the global relaxation structure of the vacuum. The Viscous Cosmic Horizon ($R_{\mathrm{U}}$) defines the maximum coherent relaxation length of this background field. A surface-tension contrast of order $\tau_{\mathrm{D}}$ distributed across $R_{\mathrm{U}}$ produces the characteristic restorative stress
\begin{equation}
    \left| \frac{\mathrm{d}\tau}{\mathrm{d}R} \right| \approx \frac{\tau_{\mathrm{D}}}{R_{\mathrm{U}}}.
\end{equation}

The stress acting on the lattice surface inertia defines the Lattice Impedance Acceleration ($a_{\mathrm{D}}$):
\begin{equation}
    \frac{\tau_{\mathrm{D}}}{R_{\mathrm{U}}} = \mu_{\mathrm{D}}a_{\mathrm{D}}.
\end{equation}

Substitution of
\begin{equation}
    \tau_{\mathrm{D}} = \mu_{\mathrm{D}}c_{\mathrm{E}}^2
\end{equation}
gives
\begin{equation}
    \frac{\mu_{\mathrm{D}}c_{\mathrm{E}}^2}{R_{\mathrm{U}}} = \mu_{\mathrm{D}}a_{\mathrm{D}}.
\end{equation}

The surface inertia cancels, yielding
\begin{equation}
    a_{\mathrm{D}} = \frac{c_{\mathrm{E}}^2}{R_{\mathrm{U}}}.
\end{equation}

The quantity $a_{\mathrm{D}}$ is therefore the acceleration scale associated with the horizon-limited relaxation capacity of the active FCC lattice.

\section{Connection to Vacuum Damping and the Galactic Acceleration Scale}

In \hyperref[app:vacuumdamping_cosmicredshift]{Appendix \ref*{app:vacuumdamping_cosmicredshift}}, we derived the Vacuum Lattice Damping Rate ($\Gamma_{\mathrm{D}} = c_{\mathrm{E}} / R_{\mathrm{U}}$), which replaces the empirical Hubble constant. We can express the Lattice Impedance directly in terms of this viscosity:
\begin{equation}
    a_{\mathrm{D}} = c_{\mathrm{E}} \cdot \Gamma_{\mathrm{D}}
\end{equation}

This relation connects the characteristic vacuum relaxation rate to its corresponding acceleration scale. The galactic condensate response and the propagation damping mechanism share the same lattice relaxation scale, although they represent different mechanical channels. Propagation damping is dissipative, whereas the coherent galactic response is reactive.

Using
\begin{equation}
    R_{\mathrm{U}} \approx 1.2665 \times 10^{26} \text{ m},
\end{equation}
the fundamental impedance acceleration is
\begin{equation}
    a_{\mathrm{D}} = \frac{(299,792,458)^2}{1.2665 \times 10^{26}} \approx 7.0964 \times 10^{-10} \text{ m/s}^2.
\end{equation}

For the reduced circular-orbit response, TFT defines the galactic acceleration scale by the closed-cycle projection
\begin{equation}
    a_{\mathrm{gal}} \equiv \frac{a_{\mathrm{D}}}{2\pi}.
\end{equation}

Numerically,
\begin{equation}
    a_{\mathrm{gal}} \approx 1.1294 \times 10^{-10} \text{ m/s}^2.
\end{equation}

The factor $1/(2\pi)$ is the closed-azimuthal normalization adopted for the galactic constitutive response. It is distinct from the source-dependent geometric transmission factor ($\Pi(R)$).

The projected scale $a_{\mathrm{gal}}$ provides the TFT mechanical counterpart of the empirical low-acceleration scale associated with galactic rotation dynamics.

\section{The Mechanism: Inductive Reactance and Apparent Dynamical Mass}

A baryonic tension gradient displaces and phase-loads the surrounding Vacuum Interstitial Condensate. The condensate responds through a correlated restorative deformation that contributes to the local acceleration field.

This behavior is analogous to inductive reactance or added-mass loading in a continuous medium. The surrounding medium contributes to the dynamical response without becoming an independent localized matter component. Because the response is reactive, it modifies the local acceleration field without requiring a continuous loss of orbital energy.

Let $g_b(R)$ denote the spherical-equivalent baryonic acceleration. The geometry-corrected radial field defined in \hyperref[app:radial_projection]{Appendix~\ref*{app:radial_projection}} is
\begin{equation}
    g_{b,\Pi}(R) = \Pi(R)g_b(R).
\end{equation}

At the reduced constitutive level, the relative response of the Vacuum Interstitial Condensate is represented by the dimensionless coefficient
\begin{equation}
    \chi_{\mathrm{g}}(R) = \sqrt{ \frac{a_{\mathrm{gal}}}{g_{b,\Pi}(R)} }.
\end{equation}

The coefficient $\chi_{\mathrm{g}}$ acts as an effective nonlinear gravitational susceptibility of the vacuum and measures the induced condensate response relative to the directly transmitted baryonic acceleration.

In the high-acceleration regime,
\begin{equation}
    g_{b,\Pi} \gg a_{\mathrm{gal}},
\end{equation}
the response coefficient satisfies
\begin{equation}
    \chi_{\mathrm{g}} \ll 1,
\end{equation}
and the induced condensate contribution is small relative to the baryonic field.

In the low-acceleration regime,
\begin{equation}
    g_{b,\Pi} \ll a_{\mathrm{gal}},
\end{equation}
the response coefficient satisfies
\begin{equation}
    \chi_{\mathrm{g}} \gg 1,
\end{equation}
and the coherent condensate contribution becomes dynamically significant.

In the spherical and unscreened limit, where
\begin{equation}
    \Pi(R) = 1,
\end{equation}
the resulting response may be represented through an apparent effective mass,
\begin{equation}
    M_{\mathrm{eff}}(R) = M_b \left[ 1+\chi_{\mathrm{g}}(R) \right].
\end{equation}

This quantity is not an additional material mass. It is the dynamical mass that would be inferred if the condensate-supported acceleration were represented entirely through an inverse-square gravitational law. For anisotropic systems, the acceleration relation provides the more general description.

\section{The Unscreened Acceleration Law}

Before finite-coherence screening is introduced, the environmental acceleration is
\begin{equation}
    g_{\mathrm{obs}}(R) = g_{b,\Pi}(R) \left[ 1+\chi_{\mathrm{g}}(R) \right].
\end{equation}

Substitution of the condensate-response coefficient gives
\begin{equation}
    g_{\mathrm{obs}}(R) = g_{b,\Pi}(R) + \sqrt{ a_{\mathrm{gal}}g_{b,\Pi}(R) }.
\end{equation}

Using the geometric transmission relation
\begin{equation}
    g_{b,\Pi}(R) = \Pi(R)g_b(R),
\end{equation}
the acceleration law becomes
\begin{equation}
    g_{\mathrm{obs}}(R) = \Pi(R)g_b(R) + \sqrt{ a_{\mathrm{gal}}\Pi(R)g_b(R) }.
\end{equation}

The first term is the directly transmitted baryonic restorative acceleration. The second term is the induced response of the Vacuum Interstitial Condensate.

In the coherent low-acceleration regime,
\begin{equation}
    g_{b,\Pi} \ll a_{\mathrm{gal}},
\end{equation}
the direct baryonic term becomes subdominant, yielding
\begin{equation}
    g_{\mathrm{obs}}(R) \approx \sqrt{ a_{\mathrm{gal}}g_{b,\Pi}(R) }.
\end{equation}

\section{Flat Rotation Curves and the Baryonic Tully--Fisher Relation}

For a finite baryonic distribution, the enclosed mass approaches an approximately constant value ($M_b$) outside the principal matter region. The spherical-equivalent baryonic acceleration is then
\begin{equation}
    g_b(R) = \frac{GM_b}{R^2}.
\end{equation}

Including geometric radial transmission gives
\begin{equation}
    g_{b,\Pi}(R) = \Pi(R) \frac{GM_b}{R^2}.
\end{equation}

In the coherent low-acceleration regime,
\begin{equation}
    g_{\mathrm{obs}}(R) \approx \sqrt{ a_{\mathrm{gal}} \Pi(R) \frac{GM_b}{R^2} }.
\end{equation}

Therefore,
\begin{equation}
    g_{\mathrm{obs}}(R) \approx \frac{ \sqrt{ a_{\mathrm{gal}}\Pi(R)GM_b } }{R}.
\end{equation}

For circular motion,
\begin{equation}
    g_{\mathrm{obs}}(R) = \frac{v^2(R)}{R}.
\end{equation}

Substitution gives
\begin{equation}
    v^2(R) = \sqrt{ a_{\mathrm{gal}}\Pi(R)GM_b }.
\end{equation}

Squaring yields
\begin{equation}
    v^4(R) = \Pi(R) G a_{\mathrm{gal}} M_b.
\end{equation}

In the outer region, the bounded geometric transmission factor approaches the spherical far-field limit,
\begin{equation}
    \Pi(R) \rightarrow 1.
\end{equation}

The asymptotic relation becomes
\begin{equation}
    v_\infty^4 = G a_{\mathrm{gal}} M_b.
\end{equation}

The orbital speed is therefore independent of radius in the coherent outer regime, producing a flat rotation curve and the Baryonic Tully--Fisher scaling.

\section{The Unified Screening Function}

The coherent condensate response cannot persist over all accelerations and distances. It is bounded by an inner activation condition and an outer coherence condition.

In high-acceleration regions, local deformation is accommodated without establishing an extended condensate response. At sufficiently large radius, the source-correlated phase relation becomes indistinguishable from the background vacuum fluctuations, terminating the galactic coupling.

These two limits are represented by the screening function
\begin{equation}
    \Xi(R) = \left[ 1+ \left( \frac{g_{b,\Pi}(R)}{g_*} \right)^{n_1} \right]^{-1} \left[ 1+ \left( \frac{R}{R_{\mathrm{c}}} \right)^{n_2} \right]^{-1}.
\end{equation}

The inner threshold ($g_*$) defines the acceleration above which coherent condensate loading is suppressed. A provisional galactic range is
\begin{equation}
    g_* \sim (50\text{--}100) a_{\mathrm{gal}}.
\end{equation}

The coherence radius ($R_{\mathrm{c}}$) defines the outer extent of the source-correlated condensate response. A representative range is
\begin{equation}
    R_{\mathrm{c}} \sim 200\text{--}500 \text{ kpc}.
\end{equation}

The positive exponents $n_1$ and $n_2$ control the sharpness of the inner and outer transitions. The baseline smooth-transition model uses
\begin{equation}
    n_1 \approx n_2 \approx 2.
\end{equation}

With finite-coherence screening, the global acceleration law becomes
\begin{equation}
    g_{\mathrm{obs}}(R) = g_{b,\Pi}(R) + \Xi(R) \sqrt{ a_{\mathrm{gal}}g_{b,\Pi}(R) }.
\end{equation}

Equivalently,
\begin{equation}
    g_{\mathrm{obs}}(R) = \Pi(R)g_b(R) + \Xi(R) \sqrt{ a_{\mathrm{gal}}\Pi(R)g_b(R) }.
\end{equation}

The screening function acts only on the induced condensate response. The directly transmitted baryonic term remains present in every regime.

\section{Limiting Behavior}

\paragraph{Inner high-acceleration region}
For
\begin{equation}
    g_{b,\Pi} \gg g_*,
\end{equation}
the inner screening factor approaches zero:
\begin{equation}
    \Xi(R) \rightarrow 0.
\end{equation}
The acceleration reduces to
\begin{equation}
    g_{\mathrm{obs}}(R) \rightarrow g_{b,\Pi}(R).
\end{equation}

\paragraph{Coherent intermediate region}
For
\begin{equation}
    g_{b,\Pi} \lesssim g_*, \qquad R \ll R_{\mathrm{c}},
\end{equation}
the screening function approaches unity:
\begin{equation}
    \Xi(R) \approx 1.
\end{equation}
The environmental acceleration becomes
\begin{equation}
    g_{\mathrm{obs}}(R) \approx g_{b,\Pi}(R) + \sqrt{ a_{\mathrm{gal}}g_{b,\Pi}(R) }.
\end{equation}

\paragraph{Deep galactic regime}
For
\begin{equation}
    g_{b,\Pi} \ll a_{\mathrm{gal}}, \qquad R \ll R_{\mathrm{c}},
\end{equation}
the condensate contribution dominates:
\begin{equation}
    g_{\mathrm{obs}}(R) \approx \sqrt{ a_{\mathrm{gal}}g_{b,\Pi}(R) }.
\end{equation}
Where $\Pi(R) \rightarrow 1$, this produces
\begin{equation}
    v_\infty^4 = G a_{\mathrm{gal}} M_b.
\end{equation}

\paragraph{Beyond the coherence radius}
For
\begin{equation}
    R \gg R_{\mathrm{c}},
\end{equation}
the outer screening factor approaches zero:
\begin{equation}
    \Xi(R) \rightarrow 0.
\end{equation}
The acceleration returns to the directly transmitted baryonic field:
\begin{equation}
    g_{\mathrm{obs}}(R) \rightarrow g_{b,\Pi}(R).
\end{equation}

The outer cutoff prevents the $1/R$ impedance regime from extending indefinitely. Without this termination, the spherical-equivalent dynamical mass would increase linearly with radius and the corresponding potential would grow logarithmically without bound.

\section{Scope and Validation}

The cosmic-impedance law is a reduced galaxy-scale constitutive model for the coherent response of the Vacuum Interstitial Condensate. The fundamental acceleration ($a_{\mathrm{D}}$) is fixed by the Absolute Optical Horizon and the vacuum relaxation rate, while ($a_{\mathrm{gal}}$) is the closed-cycle scale entering the galactic response.

The factor $\Pi(R)$ describes source-dependent geometric transmission. The function $\Xi(R)$ controls the inner activation and outer coherence of the induced condensate contribution. The parameters $g_*$, $R_{\mathrm{c}}$, $n_1$, and $n_2$ require calibration against galaxy rotation curves, high-acceleration orbital constraints, weak-lensing profiles, and numerical solutions of the coupled baryonic--condensate field.

The present formulation addresses the galaxy-scale excess-acceleration regime. Applications to galaxy clusters, merging systems, and nonspherical lensing geometries require the corresponding multidimensional tension and condensate-coherence solutions.
